\section{Фиксированная размерность}\label{app:fixed-dim}

Нижние оценки раздела~\ref{sec:lower} используют растущую размерность. Покажем, что это существенно: в фиксированной размерности метод центров тяжести решает задачу при фиксированной нормированной точности $L_1^2\varepsilon/L_0$ за число обращений, логарифмическое по $L_1R$.

Функция $t\mapsto e^t-1-t$ строго возрастает на $[0,\infty)$ от $0$ до $+\infty$. Поэтому однозначно определено число $\rho_\varepsilon>0$ --- радиус шара вокруг точки минимума, в котором оценка роста леммы~\ref{lem:radial} не превосходит $\varepsilon$:
\begin{equation}\label{eq:fixed-radius}
 \frac{L_0}{L_1^2}\bigl(e^{L_1\rho_\varepsilon}-1-L_1\rho_\varepsilon\bigr)=\varepsilon .
\end{equation}

\begin{proposition}[метод центров тяжести]\label{prop:fixed-dim}
Пусть $L_0,L_1,R,\varepsilon>0$ и размерность $d$ фиксирована. Класс $\mathcal P_{\rm R}$ в размерности $d$ решается методом центров тяжести за
\begin{equation}\label{eq:centroid-count}
 1+\Bigl\lfloor\frac{d\log\bigl(2\sqrt d\,R/\min\{R,\rho_\varepsilon\}\bigr)}{-\log(1-1/e)}\Bigr\rfloor
\end{equation}
обращений к оракулу, то есть за
\[
 O\Bigl(d\log\Bigl(2\sqrt d\max\Bigl\{1,\frac R{\rho_\varepsilon}\Bigr\}\Bigr)\Bigr)
\]
обращений. В размерности $d=1$ метод деления отрезка пополам использует не более
\begin{equation}\label{eq:one-dim-count}
 \Bigl\lceil\max\Bigl\{\log_2\frac R{\rho_\varepsilon},\,0\Bigr\}\Bigr\rceil
\end{equation}
обращений. Если выход обязан совпадать с точкой обращения к оракулу, достаточно добавить одно обращение.
\end{proposition}

При фиксированных $d$ и $L_1^2\varepsilon/L_0$ величина $L_1\rho_\varepsilon$ постоянна, и оценка~\eqref{eq:centroid-count} имеет порядок $\log(2+L_1R)$. Многомерный метод центров тяжести выбирает выход по значениям функции. Метод деления отрезка пополам использует только знаки производной.

\begin{proof}
\leavevmode
\begin{itemize}
\item \emph{Шар точек с малой невязкой.} По лемме~\ref{lem:radial} для каждого $x$
\begin{equation}\label{eq:fixed-envelope}
 f(x)-f^\ast\le\frac{L_0}{L_1^2}\bigl(e^{L_1\norm{x-x^\ast}}-1-L_1\norm{x-x^\ast}\bigr).
\end{equation}
Следовательно, в шаре $B(x^\ast,\min\{R,\rho_\varepsilon\})$ невязка по функции не больше $\varepsilon$.
\item \emph{Метод центров тяжести.} Без ограничения общности $x^0=0$. Начнём с куба $K_0=[-2R,2R]^d$. Он содержит шар $B(x^\ast,\min\{R,\rho_\varepsilon\})$, поскольку $\norm{x^\ast}\le R$. На шаге $k$ вычислим центр тяжести $c_k$ многогранника $K_{k-1}$ и обратимся к оракулу в этой точке. Если $\nabla f(c_k)=0$, то $c_k$ --- точка минимума. Иначе положим
\[
 K_k=K_{k-1}\cap\bigl\{z:\ \ip{\nabla f(c_k)}{z-c_k}\le0\bigr\}.
\]
После заданного числа шагов вернём точку обращения с наименьшим значением функции. Вычисления центров тяжести между обращениями к оракулу не учитываются.
\item \emph{Шар сохраняется.} Пусть во всех первых $N$ точках обращения невязка по функции больше $\varepsilon$. Для всякой точки $z$ шара $B(x^\ast,\min\{R,\rho_\varepsilon\})$ из выпуклости и~\eqref{eq:fixed-envelope}
\[
 \ip{\nabla f(c_k)}{z-c_k}\le f(z)-f(c_k)<0 .
\]
Значит, весь шар сохраняется после каждого отсечения.
\item \emph{Сравнение объёмов.} По теореме Грюнбаума \citep{grunbaum1960partitions} всякая гиперплоскость, проходящая через центр тяжести выпуклого тела, оставляет в каждом из двух полупространств не менее $1/e$ его объёма. Единичный шар содержит куб $[-1/\sqrt d,1/\sqrt d]^d$, поэтому его объём не меньше $(2/\sqrt d)^d$. Следовательно,
\begin{align*}
 \Bigl(\frac{2\min\{R,\rho_\varepsilon\}}{\sqrt d}\Bigr)^d
 &\by{\le}{(\mathrm a)}\operatorname{vol}(K_N)\\
 &\by{\le}{(\mathrm б)}\Bigl(1-\frac1e\Bigr)^N(4R)^d ,
\end{align*}
где (а) --- сохранение шара, (б) --- теорема Грюнбаума.
При $N$, равном~\eqref{eq:centroid-count}, правая часть строго меньше левой. Это противоречие показывает, что хотя бы одна из точек обращения к оракулу имеет невязку не больше $\varepsilon$.
\item \emph{Одномерный случай.} Начнём с отрезка $[-R,R]$. На каждом шаге обращаемся к оракулу в середине текущего отрезка. Если производная в ней положительна, оставляем левую половину, если отрицательна --- правую, а при нулевой производной найден точный минимум. Монотонность производной выпуклой функции гарантирует, что точка минимума остаётся в отрезке. После $n$ обращений середина $\hat x$ оставшегося отрезка удовлетворяет неравенствам
\begin{align*}
 |\hat x-x^\ast|&\bypad{\le}{}{\eqref{eq:fixed-envelope}}R\,2^{-n},\\
 f(\hat x)-f^\ast&\by{\le}{\eqref{eq:fixed-envelope}}\frac{L_0}{L_1^2}\bigl(e^{L_1R2^{-n}}-1-L_1R\,2^{-n}\bigr).
\end{align*}
При $n$, равном~\eqref{eq:one-dim-count}, имеем $R2^{-n}\le\rho_\varepsilon$, и правая часть не больше $\varepsilon$. В выходной середине к оракулу могли не обращаться. Если выход должен совпадать с точкой обращения, нужно ещё одно обращение.\qedhere
\end{itemize}
\end{proof}

\begin{remark}\label{rem:fixed-dim-growth}
При фиксированных $d$ и $L_1^2\varepsilon/L_0$ предложение~\ref{prop:fixed-dim} даёт лишь $O(\log(2+L_1R))$ обращений. Поэтому нижняя оценка, линейная по $L_1R$ с точностью до логарифмов, не может быть реализована семейством фиксированной размерности. Ей необходим механизм сокрытия информации в числе направлений, растущем вместе с числом обращений. Именно поэтому в теоремах~\ref{thm:lower-gap} и~\ref{thm:lower-radius} используется размерность порядка длины цепочки.
\end{remark}
